\section{Related Work}

Work related to PRISM falls into four broad categories: prompt-injection and jailbreak defense, security mechanisms for tool-using agents, runtime monitoring and policy enforcement, and tamper-evident audit mechanisms. PRISM is closest to work that treats agent security as a runtime systems problem rather than as a standalone text-classification task. At the same time, PRISM does not claim to be the first secure-agent framework or the first guardrail system. Its distinction is narrower and more concrete: a zero-fork runtime security layer for OpenClaw-based agent gateways that combines lifecycle-wide interception, staged risk response, policy-enforced tool and network controls, and an operational audit plane within a single deployable system.

\subsection{Prompt-Injection and Boundary Guardrails}

One important line of work studies prompt injection and jailbreak defense primarily at the level of text classification or boundary filtering. Early work by Greshake et al.~\cite{greshake2023} established indirect prompt injection as a concrete threat to LLM-integrated applications by showing how adversaries can plant malicious instructions in retrieved content rather than in direct user prompts. Survey-style work such as \emph{Security of AI Agents}~\cite{securityofaiagents2024} provides a broader view of the attack surface introduced by agentic systems---including prompt override, unsafe tool use, and data leakage---and helps motivate why prompt-only defenses are insufficient in agent settings. Benchmark-oriented work such as \emph{Agent Security Bench (ASB)}~\cite{asb2024} and large-scale prompt-hacking studies such as \emph{HackAPrompt}~\cite{hackaprompt2023} further highlight that attacks can manifest across multiple stages of agent operation rather than at a single input boundary.

More targeted defenses have explored firewall-style protection against indirect prompt injection. \emph{Indirect Prompt Injections: Are Firewalls All You Need, or Stronger Benchmarks?}~\cite{indirectpromptfirewalls2025}, for example, frames defense in terms of a tool-input firewall (Minimizer) and a tool-output firewall (Sanitizer) placed around untrusted content flows. PRISM shares the same concern that indirect injection often enters through fetched content or tool-returned text, but differs in scope: rather than treating tool input and tool output as the only security boundaries, PRISM distributes enforcement across message ingress, prompt construction, pre-execution tool checks, post-tool scanning, persistence-time redaction, outbound filtering, and session-lifecycle hooks. The relevant comparison is therefore not that PRISM replaces boundary firewalls, but that it generalizes boundary filtering into a broader runtime-defense model for tool-augmented gateways.

\subsection{Security Mechanisms for Tool-Using Agents}

A second line of work focuses on guardrails designed specifically for LLM agents that use tools. \emph{LlamaFirewall}~\cite{llamafirewall2025} presents an open-source guardrail framework for secure AI agents and represents one of the closest points of comparison to PRISM. \emph{NeMo Guardrails}~\cite{nemoguardrails2023} likewise provides programmable runtime rails for controllable and safe LLM applications. Both systems treat agent security as an operational problem rather than a single detection benchmark. The primary difference lies in framing and integration depth: these systems are presented as more general guardrail frameworks, whereas PRISM is deliberately implemented as a zero-fork runtime layer attached to one concrete gateway stack. This narrower scope allows PRISM to lean more heavily on gateway lifecycle hooks, deploy-time policy enforcement, and operator workflows, at the cost of reduced framework generality.

Other work employs dedicated guard agents or adaptive safety mechanisms to supervise agent behavior. \emph{GuardAgent}~\cite{guardagent2024} proposes a guard-agent approach based on knowledge-enabled reasoning, while \emph{AGrail}~\cite{agrail2025} emphasizes adaptive and lifelong safety checks for agent tasks. These approaches are complementary to PRISM but reflect different design philosophies. Their primary emphasis is on supervisory reasoning or adaptive safety detection; PRISM instead emphasizes explicit interception points, policy-backed runtime controls, and recoverable operational behavior. PRISM is not built around the assumption that a single supervisory model can reliably arbitrate all unsafe behavior; it is built around the premise that agent security benefits from multiple concrete enforcement points tied directly to the runtime itself.

Recent work on safe tool use makes a distinct but highly relevant contribution. \emph{Towards Verifiably Safe Tool Use for LLM Agents}~\cite{safe-tool-use2026} argues for deriving enforceable safety requirements from structured workflow analysis and formalized tool constraints. This direction is more ambitious in terms of formal assurance than PRISM currently pursues. PRISM should therefore be understood as complementary rather than competing on the same axis: it provides deployable runtime controls and auditing for an existing gateway implementation, while formal safe-tool-use methods point toward a stronger future direction for specifying and verifying tool-level safety properties.

\subsection{Runtime Monitoring, Policy Enforcement, and Benchmarks}

A third category encompasses work that treats agent security as a runtime-control or evaluation problem. Some of this literature focuses on formalizing attacks and defenses, while other work focuses on adaptive guardrails or workflow-level constraints. \emph{Agent Security Bench (ASB)}~\cite{asb2024} is particularly relevant here because it frames agent security across multiple scenarios, tools, and attack types, demonstrating that vulnerabilities span prompt handling, tool usage, and memory-related stages. \emph{ToolEmu}~\cite{toolemu2023} is also relevant because it provides a way to evaluate tool-using agents against high-stakes scenarios without fully instantiating each external tool environment. PRISM is not a competing benchmark; rather, it is a deployable defense architecture that could in principle be evaluated using benchmark methodologies of this kind.

The most important conceptual distinction in this category is between detector-style guardrails and policy-enforced runtime defense. Detector-style systems primarily answer whether a prompt or tool result appears suspicious. PRISM includes detection as well, but it couples detection with specific enforcement controls: high-risk tool blocking, protected-path enforcement, private-network restrictions, tiered-domain handling, outbound secret filtering, and operator-mediated exception workflows. This coupling is significant because many agent failures are not best characterized as classification errors alone---they are failures of execution policy, state management, or unsafe privilege exposure. In that sense, PRISM is closer to a runtime security layer than to a pure prompt-defense classifier.

\subsection{Tamper-Evident Auditing and Operational Security}

A fourth related line of work originates from the secure-logging and tamper-evident auditing literature. Prior systems outside the agent-security domain have long studied append-only logs, chained integrity records, and post-hoc verification mechanisms for detecting unauthorized modification of historical events; Schneier and Kelsey~\cite{schneierkelsey1998}, for example, provide an early cryptographic formulation of secure logging on untrusted machines. PRISM does not claim a new cryptographic audit primitive. Instead, it adapts established tamper-evident techniques to an agent-runtime setting by combining HMAC-protected audit entries, chained hashes, anchor-based verification, and CLI verification tooling with a runtime dashboard and hot-reloadable policy plane.

This operational coupling is an important distinction from detector-only guardrail systems. In PRISM, auditability is not a secondary observability feature added after detection; it is integral to the security argument for recoverability, operator review, and policy-change management. The same holds for health probes, configuration reloads, and allow workflows. These features are not algorithmic contributions in the model-research sense, but they help distinguish PRISM from work that primarily studies detection quality while leaving runtime operations outside the system boundary.

Overall, PRISM sits at the intersection of these four threads: it is broader than boundary-only prompt firewalls, more operationally grounded than detector-only guard agents, less formal than specification-driven safe-tool-use work, and more agent-specific than generic tamper-evident logging systems. The resulting contribution is not a new classifier or a new cryptographic primitive, but a deployable runtime security architecture that integrates multiple defense and operations mechanisms into a single gateway-facing system.
